% ---
% filename : ["electrotechnique_circuits_ac_20260428_445_bobine_impedance_inductance.tex"]
% id : "electrotechnique_circuits_ac_20260428_445"
% title : "Bobine, impédance et inductance"
% domain : "Électrotechnique"
% subdomain : "Circuits AC"
% tags : ["bobine", "résistance", "impédance", "inductance", "réactance", "courant alternatif", "ts"]
% difficulty : 2
% time_solve : 12
% has_octave : True
% octave_files : ["electrotechnique_circuits_ac_20260428_445_bobine_impedance_inductance.m"]
% author : "om"
% ---
\documentclass[a4paper,11pt,answers,addpoints]{exam}

% ---------------------------------------------------------
% LANGUE, ENCODAGE, FONTS
% ---------------------------------------------------------
\usepackage[utf8]{inputenc}

% ---------------------------------------------------------
% GÉOMÉTRIE ET MISE EN PAGE
% ---------------------------------------------------------
\usepackage[left=1.7cm,right=1.5cm,top=2cm,bottom=1cm]{geometry}
%\usepackage{a4wide}
\usepackage{microtype}      % Meilleure répartition du texte
\usepackage{float}          % Positionnement [H]
\usepackage{needspace}      % Saut conditionnel intelligent

% Éviter les veuves/orphelines (optionnel, à décommenter si besoin)
% \widowpenalty=10000
% \clubpenalty=10000

% Espacement vertical flexible
\raggedbottom
\setlength{\parskip}{0.5em plus 0.2em minus 0.1em}
\setlength{\parindent}{0pt}


% ---------------------------------------------------------
% MATHÉMATIQUES
% ---------------------------------------------------------
\usepackage{amsmath, amssymb, bm, esvect}

% Vecteurs
\newcommand{\V}{\overrightarrow}
\def\vect#1{\overrightarrow{\strut#1}}
\providecommand{\Degrees}[0]{\ensuremath{^\circ}}

% ---------------------------------------------------------
% UNITÉS ET GRANDEURS (SIunitx)
% ---------------------------------------------------------
\usepackage{siunitx}
\sisetup{output-decimal-marker={,}}
\DeclareSIUnit \var {var}
\DeclareSIUnit \VA {VA}
\DeclareSIUnit \kVA {kVA}
\DeclareSIUnit[quantity-product={}] \degree{\text{\textdegree}}

% ---------------------------------------------------------
% GRAPHIQUES : TikZ + CircuitikZ + PGFPlots
% ---------------------------------------------------------
\usepackage{tikz}
\usetikzlibrary{
	arrows.meta,
	intersections,
	decorations.pathreplacing,
	positioning
}

\usepackage[european,straightvoltages,EFvoltages]{circuitikz}

\usepackage{tkz-fct}
\usepackage{pgfplots}
\pgfplotsset{compat=1.12}

% Symbole étoile-triangle personnalisé
\newcommand{\wye}{%
	\mathbin{\tikz[x=1ex,y=1ex]{%
			\draw[line width=.1ex] (0,0)--(45:1)--++(-45:1) (45:1)--++(0,1);
	}}%
}


\usepackage{sankey}

\pgfdeclarelayer{background}
\pgfdeclarelayer{foreground}
\pgfdeclarelayer{sankeydebug}
\pgfsetlayers{background,main,foreground,sankeydebug}


% ---------------------------------------------------------
% TABLEAUX
% ---------------------------------------------------------
\usepackage{tabularx}
\usepackage{tabu}

% ---------------------------------------------------------
% HYPERLIENS
% ---------------------------------------------------------
\usepackage[colorlinks=true,
menucolor=red,
breaklinks=true,
urlcolor=blue,
linkcolor=blue]{hyperref}

% ---------------------------------------------------------
% COMMANDES PERSONNELLES
% ---------------------------------------------------------
\newcommand\nomauteur{bdminasmorgul@protonmail.com}
\newcommand\dateTe{28.04.2026}
\newcommand\brancheTe{TS}

% ---------------------------------------------------------
% STYLE DE L'EXERCICE
% ---------------------------------------------------------
\pagestyle{headandfoot}
\firstpageheadrule
\runningheadrule

\firstpageheader{\brancheTe\ (\nomauteur)}{Élève : }{(\dateTe) \thepage/\numpages}
\runningheader{\brancheTe\ (\nomauteur)}{Élève : }{(\dateTe) \thepage/\numpages}

% Compteur en bas de page
\newcommand{\totalpointtts}{%
	\begin{tikzpicture}[remember picture, overlay]
		\node[anchor=south west, inner sep=0pt, outer sep=0pt]
		at ([xshift=0.5cm,yshift=0.5cm]current page.south west)
		{\rotatebox{90}{Total des points : \numpoints}};
	\end{tikzpicture}
}

\firstpagefooter{\totalpointtts}{}{}
\runningfooter{\totalpointtts}{}{}

% Réglages exam
\printanswers
\addpoints
\pointsinmargin
\bracketedpoints

% ---------------------------------------------------------
% LISTINGS (code)
% ---------------------------------------------------------
\usepackage{xcolor}
\usepackage{listings}

\lstdefinestyle{octavestyle}{
	language=Matlab,
	basicstyle=\ttfamily\small,
	keywordstyle=\color{blue},
	commentstyle=\color{green!50!black},
	stringstyle=\color{red},
	stepnumber=1,
	frame=single,
	breaklines=true,
	breakatwhitespace=true,
	tabsize=4
}

% ---------------------------------------------------------
% DOCUMENT
% ---------------------------------------------------------
\begin{document}
\unframedsolutions
\pointsinmargin
\bracketedpoints
\section*{Préparation TS PQ CFC ELMO, IELE, PELE}
\begin{questions}

\question En courant continu, on mesure dans une bobine un courant de \qty{0,4}{[\ampere]} sous une tension de \qty{10}{[\volt]}. En courant alternatif, la valeur de l'intensité dans la même bobine est de \qty{2}{[\ampere]} sous une tension de \qty{110}{[\volt]}~/~\qty{50}{[\hertz]}. Calculer les valeurs suivantes~?

\begin{parts}
	\part[3] La résistance de la bobine.
	\part[3] L'impédance de la bobine.
	\part[3] L'inductance de la bobine.
\end{parts}
\vspace{0.5em}
\needspace{55\baselineskip}
\begin{solution}
	\begin{align*}
		U_{DC} &= \qty{10}{[\volt]} \quad ; \quad I_{DC} = \qty{0,4}{[\ampere]} \\[6pt]
		R &= \dfrac{U_{DC}}{I_{DC}} = \dfrac{\qty{10}{[\volt]}}{\qty{0,4}{[\ampere]}} = \qty{25}{[\ohm]} \\[6pt]
		U_{AC} &= \qty{110}{[\volt]} \quad ; \quad I_{AC} = \qty{2}{[\ampere]} \\[6pt]
		Z &= \dfrac{U_{AC}}{I_{AC}} = \dfrac{\qty{110}{[\volt]}}{\qty{2}{[\ampere]}} = \qty{55}{[\ohm]} \\[6pt]
		X_L &= \sqrt{Z^2 - R^2} = \sqrt{(\qty{55}{})^2 - (\qty{25}{})^2} \\[6pt]
		X_L &= \sqrt{\qty{3025}{} - \qty{625}{}} = \qty{48,99}{[\ohm]} \\[6pt]
		L &= \dfrac{X_L}{2 \cdot \pi \cdot f} = \dfrac{\qty{48,99}{[\ohm]}}{2 \cdot \pi \cdot \qty{50}{[\hertz]}} \\[6pt]
		L &= \qty{0,1559}{[\henry]} = \qty{155,9}{[\milli\henry]} \\[6pt]
	\end{align*}
	
	\begin{verbatim}
		% Télécharger OCTAVE depuis https://wiki.octave.org/Using_Octave et l'installer
		% Puis copier-coller le code dans OCTAVE
		
		% Données de l'exercice
		U_dc = 10;      % Tension continue [V]
		I_dc = 0.4;     % Courant continu [A]
		U_ac = 110;     % Tension alternative [V]
		I_ac = 2;       % Courant alternatif [A]
		f = 50;         % Fréquence [Hz]
		
		% a) Calcul de la résistance R
		R = U_dc / I_dc;
		
		% b) Calcul de l'impédance Z
		Z = U_ac / I_ac;
		
		% c) Calcul de la réactance XL et de l'inductance L
		Xl = sqrt(Z^2 - R^2);
		L = Xl / (2 * pi * f);
		
		% Affichage des résultats intermédiaires et finaux
		printf("Resistance de la bobine (DC) : %.2f [Ohm]\n", R);
		printf("Impedance de la bobine (AC) : %.2f [Ohm]\n", Z);
		printf("Reactance inductive XL : %.2f [Ohm]\n", Xl);
		printf("Inductance L calculee : %.4f [H] soit %.2f [mH]\n", L, L*1000);
	\end{verbatim}
\end{solution}

	
\end{questions}
\end{document}
